\paragraph{The Q-Learning-guided IG wins, but the question is why}
Against LPT and seven ILS variants in the twelve synthetic classes, \texttt{ig\_ql} wins every one of the $72$ comparisons that remain significant after Holm correction (of $96$); the $24$
non-significant ones fall in four classes in which the gap available to any method is already below $0.3\%$. It also beats the reimplemented simulated annealing \citep{LiPerkins2024} and bee
colony \citep{DavidovicEtal2012} in all classes (including with the budget rules relaxed in their favour: $36$ of $36$ tests) but loses to MULTIFIT \citep{CoffmanGareyJohnson1978}
for $n\geq75$, a one-shot construction that costs about one fiftieth of the CPU time; ten times the budget leaves it $30\times$ worse at $(300,20)$. These results, in
Supplementary Section~\ref{S-sec:f1-headline}, are what a typical study would report and read as a success of the learned policy. The rest of the section asks what the policy contributes. Optimality
is certified with the linear relaxation of the arc-flow model, which equals the best makespan found in $75.3\%$ of the synthetic instances (proving optimality) and never exceeded the optimum in $150$
published instances (Supplementary Section~\ref{S-sec:lbaf}).

\paragraph{The policy adds hundredths of a point; the mechanism, tenths}
Against its four ablation siblings (random selection, UCB1, UCB1 per state, Q-Learning without state) \texttt{ig\_ql} has $7$ of $48$ raw $p$-values below $0.05$ and none after Holm correction: per
class, single-run tests cannot resolve the effect. Ten algorithm seeds on $240$ instances show that it is real but strictly negligible ($\bar d=-0.011\%$, 90\% interval $-0.019$ to $-0.003$), and
pooled tests on published instances (Section~\ref{sec:external}) put it at $0.01$--$0.03$ points. The search mechanism, by contrast, is large: with learning removed from both sides,
\texttt{ig\_random} beats \texttt{ils\_random} in $10$ of $12$ classes (the exceptions are the near-ceiling classes $(20,4)$ and $(40,5)$), by $0.49\%$ on seed averages, against $0.53\%$ for the learned versions.

\paragraph{Which part of the search structure? Equal CPU time decides}
The advantage of the IG mixes the destruction--reconstruction cycle with the absence of an exhaustive local search. A hybrid that adds a short best-improvement descent after each reconstruction separates
them (Tables~\ref{tab:f25_igls_p} and~\ref{tab:f28_igls_eq}). The pure IG already beats the ILS by $0.59\%$ ($381$ wins against $61$), the descent adds a further $0.15\%$ ($436$ against $16$) and Q-Learning adds
nothing to the hybrid ($-0.004\%$). That $0.15\%$, however, is a comparison at an equal number of iterations, and the descent multiplies the CPU time per iteration by $2.1$ to $52.4$. Giving the plain IG the
CPU time of the hybrid (calibrated sequentially on an idle machine) reduces the advantage of the local search to $0.010\%$ (90\% interval $0.002$--$0.018$, inside the $\pm0.05\%$ margin):
\textbf{$93\%$ of the gain was budget}. The active ingredient for identical machines is the destruction--reconstruction cycle and the number of iterations it receives; the descent mostly substitutes for iterations.

% Auto-generated by analysis/make_tables_cycle3.py (Entrada 125)
\begin{table}[t]
\centering
\caption{$\mathrm{P}||C_{\max}$: is the active ingredient the destruction--reconstruction cycle or the local search? IG+LS adds a best-improvement descent (inter-swap, reinsert; two passes) after each reconstruction ($T=1000$, 12 classes; the ILS has 30 instances in the six original classes). $\bar d$ in \% of the makespan (negative: first better); Wilcoxon, Holm over three.}
\label{tab:f25_igls_p}
{\scriptsize\setlength{\tabcolsep}{3.5pt}
\begin{tabular}{@{}lrrrrr@{}}
\toprule
Comparison & $N$ & W/L/T & $\bar d$ (\%) & 90\% interval & $p_{\mathrm{Holm}}$ \\
\midrule
P1: IG+LS vs.\ IG (random selection) & 1{,}200 & 436/16/748 & -0.150 & [-0.162, -0.139] & $7\times10^{-65}$ \\
P2: IG+LS vs.\ ILS (random selection) & 780 & 459/16/305 & -0.722 & [-0.796, -0.650] & $2\times10^{-75}$ \\
P3: IG+LS, \textsc{QL} vs.\ random selection & 1{,}200 & 98/77/1025 & -0.004 & [-0.011, +0.004] & 0.646 \\
\bottomrule
\end{tabular}}
\end{table}

% Auto-generated by analysis/make_tables_cycle5.py (Entrada 129)
\begin{table}[t]
\centering
\caption{Is the gain of the local search in the hybrid of Table~\ref{tab:f25_igls_p} an effect of search structure or of budget? \textsc{IG} (random selection, $d=0.3n$, no learning) with $T_{\mathrm{eq}}$ iterations has the same CPU time as the hybrid with $T=1000$ ($T_{\mathrm{eq}}=1000K$, $K$ the sequential CPU ratio per class, from $2.1$ to $52.4$; checked again on an idle machine: ratios $0.96$--$1.02$). 12 classes $\times$ 100 instances; $\bar d$ in \% of the makespan (negative: first better); Wilcoxon. Q1 was the pre-specified test; Q2 and the first row are descriptive.}
\label{tab:f28_igls_eq}
{\scriptsize\setlength{\tabcolsep}{3.5pt}
\begin{tabular}{@{}lrrrrr@{}}
\toprule
Comparison & $N$ & W/L/T & $\bar d$ (\%) & 90\% interval & $p$ \\
\midrule
IG+LS vs.\ IG at equal iterations ($T=1000$) & 1{,}200 & 436/16/748 & -0.150 & -- & -- \\
Q1: IG+LS ($T=1000$) vs.\ IG at \emph{equal CPU} ($T_{\mathrm{eq}}$) & 1{,}200 & 112/104/984 & -0.010 & [-0.018, -0.002] & 0.015 \\
Q2: IG at equal CPU vs.\ IG ($T=1000$) & 1{,}200 & 421/0/779 & -0.140 & [-0.150, -0.129] & $1\times10^{-70}$ \\
\midrule
\multicolumn{6}{l}{Certified bound (1{,}200 instances): provably optimal / mean GAP (\%)} \\
\multicolumn{6}{l}{\quad IG ($T=1000$) 36.9\% / 0.347; IG+LS ($T=1000$) 55.2\% / 0.196; IG at equal CPU 54.5\% / 0.206} \\
\bottomrule
\end{tabular}}
\end{table}


\paragraph{Why the policy adds so little}
Four explanations were tested with analyses specified beforehand (details in Supplementary Section~\ref{S-sec:m-robust}). \emph{Nothing to learn:} the best fixed destruction strategy per class is no better than random
selection and an oracle that knew the best strategy of every instance would gain only $0.037\%$. \emph{State:} the original state is nearly degenerate (two features constant), but two recalibrated
states leave \texttt{ig\_ql} indistinguishable from random destruction (pooled $p_{\mathrm{Holm}}=0.28$) and make the ILS agent significantly \emph{worse} than random operator selection ($+4.8\%$ of the GAP,
$p_{\mathrm{Holm}}=0.009$). \emph{Signal:} a dense reward makes the agent worse ($+0.084\%$), pre-training the table per class on $50$ instances leaves it equivalent to random ($+0.002\%$) and
the learning rate, discount and exploration decay change nothing ($|\bar d|\leq0.023\%$). \emph{Budget:} tightening $T$ from $1000$ to $100$ almost triples the headroom ($0.037\%\to0.103\%$) but leaves the
policy effect unchanged ($-0.011\%$ at both).

\paragraph{An action space with real headroom: the destruction size}
If the three strategies are nearly interchangeable, the destruction size $d/n\in\{0.05,\ldots,0.40\}$ is not. Table~\ref{tab:f15_dsel} reports the pre-specified test ($1{,}200$ instances). The fraction
chosen \emph{by validation} on $30$ separate instances per class beats random choice out of sample ($\bar d=-0.089\%$), so the headroom is real. Q-Learning over the same five actions does not capture it
($+0.015\%$, $2\%$ of the headroom), neither does UCB1 ($+0.007\%$), and Q-Learning is worse than the validated fraction by $0.103\%$ ($318$ of $375$ decisive instances). With five times the budget it still
captures nothing significant ($-0.010\%$, interval $-0.028$ to $+0.008$, while $0.049\%$ of headroom remains), and a table pre-trained on $50$ independent instances per class loses $0.063\%$ to validation. In this problem
the gain of the action space is realised by choosing the parameter \emph{offline}, on separate instances, not by the online or cross-instance tabular learning tested here.

% Auto-generated by analysis/make_tables_cycle2.py (Entrada 122)
\begin{table}[t]
\centering
\caption{The action is the destruction size ($d/n\in\{0.05,0.10,0.20,0.30,0.40\}$, random destruction), a space with real headroom. Pre-specified tests on 1{,}200 instances ($T=1000$; Wilcoxon, Holm over four). $\bar d$: mean relative makespan difference (\%; negative: first better); W/L/T: first better/worse/tied. ``Best $d$ by validation'' is chosen per class on 30 separate instances.}
\label{tab:f15_dsel}
{\scriptsize\setlength{\tabcolsep}{3.5pt}
\begin{tabular}{@{}lrrrr@{}}
\toprule
Comparison & W/L/T & $\bar d$ (\%) & 90\% interval & $p_{\mathrm{Holm}}$ \\
\midrule
A1: best $d$ by validation vs.\ random $d$ & 324/47/829 & -0.089 & [-0.100, -0.077] & $6\times10^{-30}$ \\
A2: \textsc{IG-QL} (over $d$) vs.\ random $d$ & 229/213/758 & +0.015 & [+0.001, +0.030] & 0.123 \\
A3: UCB1 (over $d$) vs.\ random $d$ & 165/178/857 & +0.007 & [-0.005, +0.018] & 0.366 \\
A4: \textsc{IG-QL} (over $d$) vs.\ best $d$ by validation & 57/318/825 & +0.103 & [+0.092, +0.116] & $3\times10^{-35}$ \\
\bottomrule
\end{tabular}}
\end{table}

